% =======================================================
% 5. PIANIFICAZIONE
% =======================================================
\section{Pianificazione}

\subsection{Suddivisione temporale}

La pianificazione del progetto è articolata nelle seguenti fasi principali:
\begin{itemize}
\item \textbf{Candidatura\textsubscript{G} e assegnazione dell'appalto}: Marzo 2026;
\item \textbf{Fase RTB\textsubscript{G} (Requirements and Technology Baseline)}: Aprile 2026 -- Giugno 2026;
\item \textbf{Fase PB\textsubscript{G} (Product Baseline)}: Giugno 2026 -- Settembre 2026.
\end{itemize}

\subsection{Pianificazione a lungo termine}

\subsubsection{Attività di candidatura (CC)}

Prima dell'avvio formale degli sprint RTB il team ha svolto le seguenti attività
preparatorie, relative alla fase di candidatura:
\begin{itemize}
\item Primo allineamento su struttura documentale, Norme di Progetto
(2026/03/13 -- in aggiornamento continuo) e Glossario (2026/03/13 -- in aggiornamento continuo).
\item Stima dell'impegno orario complessivo e definizione del budget di progetto,
presentati nel Preventivo dei Costi\textsubscript{G} e Assunzione degli
Impegni (2026/03/13 -- 2026/03/26).
\item Analisi e valutazione della fattibilità dei capitolati proposti, con
motivazione della scelta del capitolato C6, proposto da Zucchetti S.p.A.,
e presentazione della Valutazione dei Capitolati\textsubscript{G}
(2026/03/17 -- 2026/03/26).
\item Stesura e presentazione formale della candidatura tramite Lettera di
Presentazione\textsubscript{G} (2026/03/27).
\item Configurazione iniziale del repository\textsubscript{G} documentale e del
sito di presentazione della documentazione ufficiale.
\end{itemize}

\subsubsection{Attività preparatorie alla RTB}
\label{sec:pre-rtb}

A seguito dell'aggiudicazione del capitolato, prima dell'avvio formale degli
sprint RTB, il gruppo ha svolto le seguenti attività preparatorie:
\begin{itemize}
\item Definizione dell'assetto organizzativo ispirato a Scrum: ruoli,
strumenti, convenzioni di naming e flusso di lavoro su GitHub.
\item Configurazione del board GitHub Project\textsubscript{G} e del Google
Form\textsubscript{G} per la creazione standardizzata delle issue,
con relativo foglio Google Sheets per il tracciamento delle ore.
\end{itemize}

\subsubsection{Attività fissate per la RTB}
\label{sec:attivita-rtb}

La fase RTB copre il periodo Aprile -- Giugno 2026. Di seguito è riportato lo
stato di avanzamento di ciascun documento o attività al momento dell'ultimo
aggiornamento del presente documento.

\paragraph{Norme di Progetto (NdP)}
Parti previste: introduzione, tailoring, processi primari, processi di supporto,
processi organizzativi, convenzioni redazionali, gestione della configurazione,
verifica, validazione, qualità e miglioramento.\\
Sprint di redazione: tutti gli sprint RTB.\\
\textbf{Data prima stesura:} 2026/03/13.\\
\textbf{Stato attuale:} in aggiornamento continuo.

\paragraph{Glossario}
Parti previste: termini tecnici, acronimi ed espressioni ambigue utilizzate nella
documentazione di progetto.\\
Sprint di redazione: tutti gli sprint RTB.\\
\textbf{Data prima stesura:} 2026/03/13.\\
\textbf{Stato attuale:} in aggiornamento continuo.

\paragraph{Piano di Progetto (PdP)}
Parti previste: introduzione, analisi dei rischi, modello di sviluppo,
Definition of Done\textsubscript{G}, pianificazione, preventivo, consuntivi per
sprint, preventivo a finire, organigramma e rotazione dei ruoli.\\
Sprint di redazione: tutti gli sprint RTB.\\
\textbf{Data prima stesura:} 2026/04/09.\\
\textbf{Stato attuale:} struttura completata con lo Sprint 8; la sezione relativa
agli sprint, ai consuntivi, alle risorse residue e al Preventivo a Finire viene
aggiornata al termine di ogni sprint sulla base del consuntivo di periodo,
della retrospettiva, delle misure correttive individuate e della revisione della
pianificazione futura.

\paragraph{Analisi dei Requisiti (AdR)}
Parti previste: introduzione, casi d'uso, requisiti funzionali e non funzionali,
tracciamento requisiti--fonti, casi d'uso--requisiti e fonti--requisiti,
riepilogo.\\
Sprint di redazione: Sprint 1 -- Sprint 8.\\
\textbf{Data prima stesura:} 2026/04/08.\\
\textbf{Stato attuale:} completata per RTB e validata dalla proponente in data
2026/05/26; eventuali aggiornamenti successivi vengono tracciati tramite issue
e riportati nei consuntivi di sprint.

\paragraph{Piano di Qualifica\textsubscript{G} (PdQ)}
Parti previste: introduzione, metriche di qualità di processo, metriche di
qualità di prodotto, cruscotto di valutazione, resoconto delle misurazioni e
azioni di miglioramento.\\
Sprint di redazione: Sprint 3 -- Sprint 8 per la struttura iniziale; aggiornamento progressivo fino alla RTB.\\
\textbf{Data prima stesura:} 2026/05/11.\\
\textbf{Stato attuale:} struttura completata con lo Sprint 8; metriche,
misurazioni, cruscotto e azioni di miglioramento vengono aggiornati
periodicamente secondo quanto definito nelle Norme di Progetto, sulla base
delle misurazioni raccolte e delle azioni correttive individuate.

\paragraph{Analisi dello Stato dell'Arte}
Parti previste: analisi delle tecnologie candidate per il sistema
\textit{Second Brain}, valutazione comparativa, scelte motivate e stack
tecnologico individuato.\\
Sprint di redazione: Sprint 5 -- Sprint 8.\\
\textbf{Data prima stesura:} 2026/05/13.\\
\textbf{Stato attuale:} completato e validato dalla proponente in data
2026/05/26; le tecnologie selezionate costituiscono la base per il setup e
l'implementazione del Proof of Concept.

\paragraph{Proof of Concept (PoC)}
Artefatto eseguibile per la validazione tecnologica della Technology Baseline.\\
Sprint di setup: Sprint 8.\\
Sprint di implementazione: Sprint 9 -- Sprint 10.\\
\textbf{Stato attuale:} completato.

Il setup del PoC comprende la predisposizione della struttura nella cartella
\texttt{poc/}, la separazione tra frontend e backend, la configurazione iniziale
di Vue.js, FastAPI e Docker Compose\textsubscript{G}. L'implementazione
successiva ha lo scopo di dimostrare la fattibilità tecnica delle tecnologie
selezionate e dei principali flussi funzionali previsti per la RTB.

\subsubsection{Attività fissate per la PB}
\label{sec:attivita-pb}

La fase PB copre il periodo Giugno -- Settembre 2026, a partire dal termine
della Revisione RTB (2026-06-23). Rispetto alla RTB, l'obiettivo si sposta
dalla validazione tecnologica (Technology Baseline) alla dimostrazione di un
prodotto sostanzialmente completo (Product Baseline), tramite un Minimum
Viable Product\textsubscript{G} (MVP) che copre la totalità dei requisiti
obbligatori individuati nell'Analisi dei Requisiti.

\paragraph{Norme di Progetto (NdP), Glossario, Piano di Progetto (PdP)}
Aggiornati in continuo per l'intera fase PB, analogamente a quanto avvenuto in
RTB: la NdP recepisce le convenzioni di codifica, test, CI/CD\textsubscript{G} e
code review maturate durante l'implementazione dell'MVP; il Glossario si
amplia con i termini tecnici emersi dalla progettazione e dallo sviluppo; il
Piano di Progetto (il presente documento) viene portato in versione
\textbf{v2.0.0}, con consuntivo RTB, Preventivo a Finire aggiornato e
pianificazione estesa alla PB.\\
Sprint di redazione: dallo Sprint successivo alla Revisione RTB fino alla
Revisione PB.\\
\textbf{Data prima stesura aggiornamento PB:} 2026/06/24 (Sprint 12).\\
\textbf{Stato attuale:} in aggiornamento continuo.

\paragraph{Analisi dei Requisiti (AdR) v2.x.x}
Raffinamento dei requisiti e dei casi d'uso sulla base del PoC, dei colloqui di revisione
RTB e del confronto con la proponente; tracciamento aggiornato verso Specifica
Tecnica e Piano di Qualifica.\\
Sprint di redazione: Sprint 12 -- Sprint 19.\\
\textbf{Data prima stesura aggiornamento PB:} 2026/06/24.\\
\textbf{Stato attuale:} aggiornata al 2026-08-24, coerente con l'MVP
implementato.

\paragraph{Piano di Qualifica (PdQ) v2.x.x}
Aggiornamento di metriche di prodotto e processo, Specifica dei Test (unità,
integrazione, sistema, accettazione) e cruscotto di valutazione della
qualità.\\
Sprint di redazione: Sprint 12 -- Sprint 20.\\
\textbf{Data prima stesura aggiornamento PB:} 2026/06/24.\\
\textbf{Stato attuale:} aggiornata al 2026-08-24; Test di Unità, Integrazione
e Sistema tutti implementati, integrati in CI\textsubscript{G} e superati, come anche i 
Test di Accettazione. Metriche, misurazioni, cruscotto e azioni di miglioramento vengono aggiornati
periodicamente secondo quanto definito nelle Norme di Progetto, sulla base
delle misurazioni raccolte e delle azioni correttive individuate.

\paragraph{Specifica Tecnica (ST) v1.x.x}
Descrive l'architettura logica e di deployment del prodotto, i pattern
architetturali e di progettazione adottati (Model-View-Controller\textsubscript{G}
"pull model", Strategy\textsubscript{G}, Factory\textsubscript{G},
Decorator\textsubscript{G}, Proxy\textsubscript{G}), i diagrammi UML e il
tracciamento componenti--requisiti.\\
Sprint di redazione: in parallelo all'implementazione dell'MVP (Sprint 12 -- Sprint 20).\\
\textbf{Data prima stesura:} 2026/06/24.\\
\textbf{Stato attuale:} completata e coerente con l'architettura
realizzata.

\paragraph{Implementazione del Minimum Viable Product (MVP)}
Sviluppo delle funzionalità obbligatorie previste dal capitolato: editor
Markdown\textsubscript{G} con formattazione, gestione di link, elenchi e
tabelle, Undo/Redo\textsubscript{G}, assistente AI (riassunto, traduzione,
riscrittura, Distant Writing\textsubscript{G}, metodo dei Sei Cappelli per
Pensare\textsubscript{G}), persistenza ed esportazione delle note (PDF, HTML,
JSON), corredata da pipeline di CI (analisi statica, type check, test e
copertura) e da suite di test automatizzati sui quattro livelli previsti
(unità, integrazione, sistema, accettazione).\\
Sprint di implementazione: dallo Sprint successivo alla Revisione RTB fino
alla Revisione PB.\\
\textbf{Stato attuale:} funzionalità obbligatorie implementate; Test di Unità,
Integrazione e Sistema superati (511/511 casi automatizzati); pipeline CI
verde con copertura ampiamente sopra soglia (95\% backend, $\sim$89,8\% linee
frontend).

\paragraph{Manuale Utente (MU) v1.0.0}
Guida pratica all'uso del prodotto dal punto di vista dell'utente finale,
redatta una volta stabilizzate le funzionalità dell'MVP.\\
Sprint di redazione: Sprint 18 -- Sprint 20.\\
\textbf{Data prima stesura:} 2026/08/04.\\
\textbf{Stato attuale:} completata e coerente con l'MVP
realizzato.

\paragraph{Validazione dell'MVP con la proponente}
Demo del prodotto alla proponente Zucchetti S.p.A. e raccolta
dell'approvazione formale richiesta dal regolamento di progetto (mail di
conferma o dichiarazione esplicita a verbale), propedeutica alla Revisione
PB.\\
\textbf{Stato attuale:} superata (\link{https://synergo-unit-unipd.github.io/Second-Brain/docs/PB/doc_gestionali/verb_esterni/verb_est_2026_09_01/verb_est_2026_09_01_firmato.pdf}{Verbale di validazione, 2026-09-01}).

\subsubsection{Diagramma di Gantt}

Il diagramma di Gantt mostra la dislocazione temporale delle attività del
progetto, le loro dipendenze e il parallelismo tra documenti e artefatti
sviluppati, verificati e raffinati durante l'avanzamento del progetto, dalla
fase di Candidatura fino alla Revisione PB.

Le dipendenze seguono il modello output--input: un'attività può iniziare solo
quando le attività da cui dipende hanno prodotto un output sufficiente.

\begin{landscape}
La struttura delle attività principali è la seguente:

\setganttlinklabel{s-s}{SS}
\setganttlinklabel{f-s}{FS}

\begin{figure}[H]
\centering
\resizebox{\linewidth}{!}{%
\begin{ganttchart}[
    time slot format=isodate,
    x unit=0.09cm,
    y unit chart=0.35cm,
    vgrid,
    hgrid,
    bar height=0.45,
    group height=0.45,
    milestone height=0.45,
    bar label font=\scriptsize,
    group label font=\scriptsize\bfseries,
    milestone label font=\scriptsize\bfseries,
    title label font=\scriptsize,
    bar/.append style={fill=gray!20},
    group/.append style={fill=gray!35},
    milestone/.append style={fill=black}
]{2026-03-13}{2026-09-10}

\gantttitlecalendar{year, month=name} \\

% ============================================================
% FASE CC
% ============================================================
\ganttgroup{Fase CC}{2026-03-13}{2026-03-27} \\

\ganttbar[name=prev]{Preventivo Costi e Impegni}{2026-03-13}{2026-03-26} \\
\ganttbar[name=valcap]{Valutazione Capitolati}{2026-03-17}{2026-03-26} \\
\ganttbar[name=ndp]{Norme di Progetto}{2026-03-13}{2026-06-23} \\
\ganttbar[name=gloss]{Glossario}{2026-03-13}{2026-06-23} \\
\ganttmilestone[name=lettera]{Lettera di Presentazione CC}{2026-03-27} \\

\ganttlink[link type=f-s]{valcap}{lettera}
\ganttlink[link type=f-s]{prev}{lettera}
\ganttlink[link type=s-s]{ndp}{gloss}

% ============================================================
% AGGIUDICAZIONE
% ============================================================
\ganttmilestone[name=aggiudicazione]{Aggiudicazione Capitolato}{2026-04-01} \\

% ============================================================
% FASE RTB
% ============================================================
\ganttgroup{Fase RTB}{2026-04-02}{2026-06-23} \\

\ganttbar[name=adr]{Analisi dei Requisiti}{2026-04-08}{2026-05-26} \\
\ganttbar[name=pdp]{Piano di Progetto}{2026-04-09}{2026-06-23} \\
\ganttbar[name=pdq]{Piano di Qualifica}{2026-05-11}{2026-06-23} \\
\ganttbar[name=soa]{Analisi Stato dell'Arte}{2026-05-13}{2026-05-26} \\
\ganttmilestone[name=validazione]{Validazione AdR e Stato dell'Arte}{2026-05-26} \\
\ganttbar[name=pocsetup]{Setup Proof of Concept}{2026-05-27}{2026-06-01} \\
\ganttbar[name=pocimpl]{Implementazione Proof of Concept}{2026-06-02}{2026-06-23} \\

% ============================================================
% DIPENDENZE RTB
% ============================================================
\ganttlink{aggiudicazione}{adr}
\ganttlink[link type=s-s]{ndp}{adr}

\ganttlink[link type=s-s]{ndp}{pdq}
\ganttlink[link type=s-s]{adr}{pdq}

\ganttlink[link type=s-s]{adr}{soa}

\ganttlink[link type=f-s]{adr}{validazione}
\ganttlink[link type=f-s]{soa}{validazione}
\ganttlink[link type=f-s]{validazione}{pocsetup}
\ganttlink[link type=f-s]{pocsetup}{pocimpl}

% ============================================================
% MILESTONE FINALE RTB
% ============================================================
\ganttmilestone[name=msrtb]{Revisione RTB}{2026-06-23} \\
\ganttlink{pdp}{msrtb}
\ganttlink{pdq}{msrtb}
\ganttlink{validazione}{msrtb}
\ganttlink{pocimpl}{msrtb}

% ============================================================
% FASE PB
% ============================================================
\ganttgroup{Fase PB}{2026-06-24}{2026-09-10} \\

\ganttbar[name=ndppb]{Norme di Progetto (PB)}{2026-06-24}{2026-09-10} \\
\ganttbar[name=glosspb]{Glossario (PB)}{2026-06-24}{2026-09-10} \\
\ganttbar[name=pdppb]{Piano di Progetto (PB)}{2026-06-24}{2026-09-10} \\
\ganttbar[name=adrpb]{Analisi dei Requisiti (PB)}{2026-06-24}{2026-08-18} \\
\ganttbar[name=st]{Specifica Tecnica}{2026-08-06}{2026-08-24} \\
\ganttbar[name=mvpimpl]{Implementazione MVP}{2026-08-14}{2026-08-21} \\
\ganttbar[name=pdqpb]{Piano di Qualifica (PB)}{2026-06-24}{2026-09-10} \\
\ganttbar[name=mu]{Manuale Utente}{2026-08-22}{2026-08-24} \\
\ganttmilestone[name=valmvp]{Validazione MVP (proponente)}{2026-09-01} \\

% ============================================================
% DIPENDENZE PB
% ============================================================
\ganttlink{msrtb}{ndppb}
\ganttlink{msrtb}{glosspb}
\ganttlink{msrtb}{pdppb}
\ganttlink{msrtb}{adrpb}
\ganttlink{msrtb}{st}
\ganttlink{msrtb}{mvpimpl}
\ganttlink{msrtb}{pdqpb}

\ganttlink[link type=s-s]{st}{mvpimpl}
\ganttlink[link type=f-s]{mvpimpl}{mu}
\ganttlink[link type=f-s]{mvpimpl}{valmvp}

% ============================================================
% MILESTONE FINALE PB
% ============================================================
\ganttmilestone[name=mspb]{Revisione PB}{2026-09-10} \\
\ganttlink{ndppb}{mspb}
\ganttlink{glosspb}{mspb}
\ganttlink{pdppb}{mspb}
\ganttlink{adrpb}{mspb}
\ganttlink{st}{mspb}
\ganttlink{pdqpb}{mspb}
\ganttlink{mu}{mspb}
\ganttlink{valmvp}{mspb}

\end{ganttchart}%
}
\caption{Diagramma di Gantt della fase CC, RTB e PB}
\end{figure}

\end{landscape}